(1) 氢样原子基态径向指数衰减尺度为 $a_0/Z$。规范波函数为 $\psi_{1s}(r)=(Z^3/(\pi a_0^3))^{1/2}e^{-Zr/a_0}$。积分核验 $4\pi(Z^3/(\pi a_0^3))\int_0^\infty r^2e^{-2Zr/a_0}dr=1$，因此可以直接计算径向期望。
(2) 利用 $\int_0^\infty r^ne^{-\alpha r}dr=n!/\alpha^{n+1}$，得到 $\langle r\rangle=4Z^3/a_0^3\int_0^\infty r^3e^{-2Zr/a_0}dr=3a_0/(2Z)$。类似地 $\langle1/r\rangle=4Z^3/a_0^3\int_0^\infty r e^{-2Zr/a_0}dr=Z/a_0$。
(3) 库仑势为 $V(r)=-Ze^2/(4\pi\varepsilon_0r)$，故 $\langle V\rangle=-Z^2e^2/(4\pi\varepsilon_0a_0)$。量子Coulomb基态总能量 $E_1=-Z^2e^2/(8\pi\varepsilon_0a_0)=-mZ^2e^4/[2(4\pi\varepsilon_0)^2\hbar^2]$，恰为势能期望的一半。
(4) 对均匀次幂势 $V\propto r^{-1}$，量子维里定理 $2\langle K\rangle=\langle r\,dV/dr\rangle=-\langle V\rangle$。于是 $\langle K\rangle=Z^2e^2/(8\pi\varepsilon_0a_0)=-E_1>0$，$\langle K\rangle+\langle V\rangle=E_1$，所有束缚能量均按 $Z^2$ 标度变化。
